\documentclass[a4paper,10pt]{article}
\usepackage{amssymb, amsmath}
\DeclareMathOperator{\arcsinh}{arcsinh}

\begin{document}

\title{Analysis Formelsammlung}
\author{Peter Merkert, Martin Thoma}
\date{21. Februar 2012}

\section{Grenzwerte}
\begin{table}[ht]
\begin{minipage}[b]{0.5\linewidth}\centering
\begin{align*}
  \lim_{x \to 0} \frac {\sin x}{x}  &= 1 \\
  \lim_{x \to 0} \frac {e^x - 1}{x} &= 1 \\
  \cos x &= \sum_{n = 0}^{\infty} (-1)^n \frac {x^{2n}}{(2n)!}
\end{align*}
\end{minipage}
\hspace{0.5cm}
\begin{minipage}[b]{0.5\linewidth}\centering
\begin{align*}
\cosh x = \frac {1}{2} (e^x + e^{-x}) &= \sum_{n = 0}^{\infty} \frac {x^{2n}}{(2n)!} \\
\sinh x = \frac {1}{2} (e^x - e^{-x}) &= \sum_{n = 0}^{\infty} \frac {x^{2n + 1}}{(2n + 1)!}
\end{align*}
\end{minipage}
\end{table}

\section{Zusammenhänge}
\begin{align*}
  (\cos x)^2 + (\sin x)^2 &= 1 \\
  (\cosh x)^2 - (\sinh x)^2 &= 1
\end{align*}

\end{document}
